\chapter{Conclusion and Future Work}
\label{ch:conclusion}

\section{Conclusion}
The objective of this thesis was to build and assess a flow-only turbulent topology optimization framework for internal configurations relevant to energy and thermal-management systems. The framework does not solve the energy equation and does not yet optimize conjugate heat transfer. Instead, it addresses the hydraulic prerequisite for those applications: generating low-loss internal flow paths under turbulent operating conditions. This narrower scope is deliberate, because turbulent momentum transport, wall-distance treatment, Brinkman penalization, geometry extraction, and adjoint sensitivity assembly already form a difficult optimization problem.

The work was organized as a sequence of checks. Chapter~\ref{ch:turbulence_modeling} assessed the useful range of the Spalart-Allmaras (SA) closure. Chapters~\ref{ch:fenics_implementation} and~\ref{ch:sa_verification} implemented and verified the FEniCS-SA solver. Chapters~\ref{ch:topology_optimization_methodology} and~\ref{ch:benchmark_results} embedded that solver in a density-based topology optimization loop with Brinkman penalization, filtering, projection, MMA, topology-dependent SA and wall-distance penalties, a frozen-turbulence adjoint, sensitivity checks, and independent sharp-geometry CFD validation.

\subsection*{Research Question 1}
\textbf{How accurately does the Spalart-Allmaras turbulence model predict curved and separated internal turbulent flows, and where do its physical limitations appear?}
\\ \\
The answer is conditional. SA is defensible as a practical design-model closure for attached wall-bounded and mildly curved internal flows, including the moderate-curvature VMFL048 U-bend evidence. It is not closure-independent truth for strongly curved or separated flows.

The Ansys comparisons in this thesis make that boundary concrete. The VMFL048 U-bend at $De\approx15{,}000$ remains credible, while the tighter VMFL030 pipe bend at $De\approx18{,}000$ already shows clear SA deterioration in the selected profile comparison. The backward-facing-step evidence gives the same warning for separated recovery: SA can predict plausible integral behavior, but SST $k-\omega$ follows the local velocity recovery more closely. The limiting issue is therefore curvature- and separation-driven Reynolds-stress anisotropy under a scalar eddy-viscosity closure, not a single Reynolds-number threshold.

\subsection*{Research Question 2}
 \textbf{Does the custom FEniCS implementation of the Spalart-Allmaras model reproduce Ansys Fluent predictions for the selected turbulent internal-flow benchmarks?}
 \\ \\
 Within the two-dimensional, wall-resolved scope of the thesis, the answer is yes. The solver reproduces the governing ingredients needed by the optimizer: incompressible flow coupling, SA transport, eddy-viscosity reconstruction, wall-distance evaluation, and near-wall resolution.

In the periodic channel, the fine-mesh bulk velocity differs from the Ansys reference by $2.67\%$, while the exported profile errors are $E_{\infty}=2.66\%$ and $E_{L_2}=2.64\%$. In the U-bend verification case, the pressure drop differs from Ansys by $2.79\%$, and the bend-exit profile gives $E_{\infty}=8.50\%$ and $E_{L_2}=2.75\%$. These results verify the FEniCS implementation as a sufficiently consistent forward solver for the topology optimization studies, while leaving the model-form limitations from Chapter~\ref{ch:turbulence_modeling} in place.

\subsection*{Research Question 3}
\textbf{To what extent do SA-driven topology optimizations produce final designs whose predicted performance remains consistent under independent Ansys Fluent validation with SA and SST $k-\omega$, and how do they compare with laminar-optimized designs under identical turbulent operating conditions?}
\\ \\
Chapter~\ref{ch:benchmark_results} demonstrates that the reliability of an SA-optimized design depends fundamentally on the flow physics dictated by the target geometry.

The pipe bend represents the cleanest transfer case, showing that the framework successfully produces robust, closure-independent designs in regimes dominated by attached flow and moderate curvature. Even as a provisional $70$-iteration endpoint, the extracted pipe-bend topology remains highly credible: replacing Ansys-SA by SST $k-\omega$ increases pressure drop by only about $4.5\%$ and viscous dissipation by $8.5\%$. Under identical SST $k-\omega$ operating conditions, the turbulent design decisively outperforms the laminar-optimized reference, reducing pressure drop by $72.0\%$ and viscous dissipation by $47.0\%$. Furthermore, the derivative supplied to MMA is internally verified, with coordinate finite differences giving a maximum relative discrepancy of $5.92\times10^{-5}$ and the directional Taylor test showing approximately second-order behavior.

Conversely, the U-bend and diffuser benchmarks expose the severe limitations of relying on a one-equation closure for topology optimization in separation-dominated flows. Because the quantitative results for these cases are highly sensitive to the validation model, the exact performance metrics lose their predictive value. For geometries forcing strong recirculation (such as the $180^\circ$ return of the U-bend) or adverse pressure gradients (the diffuser), the optimizer heavily tunes the layout to the SA model's specific, and often overly forgiving, prediction of flow separation and recovery. When these extracted shapes are cross-examined using SST $k-\omega$, which is notably more responsive to adverse pressure gradients, the predicted flow structures disagree and the performance deteriorates significantly.

Consequently, the advantage of a turbulent-optimized design over a laminar baseline is not absolute. In attached flows, the turbulent formulation is a clear necessity and delivers massive performance gains. But in heavily separated regimes, the optimized geometry becomes so deeply coupled to the SA model's simplifying assumptions that the turbulent designs offer only inconclusive, marginal gains over a simple laminar surrogate, or can even perform worse when evaluated under a different turbulence closure.

\subsection*{Overall Conclusion}
The central conclusion is balanced. A frozen-SA density-based topology optimization framework can produce credible low-loss turbulent internal-flow design candidates when the flow remains within the closure's practical range and when the extracted geometry is independently validated. It is not a substitute for final CFD validation, experimental testing, or a higher-fidelity turbulence model in separated, strongly three-dimensional flows. Independent sharp-geometry CFD validation is the mandatory engineering step between an optimized density field and anything that could be manufactured.

Within its validated envelope, the framework can still be valuable as an early-stage design and research tool for internal hydraulic layouts where pressure loss or dissipated mechanical energy is the main design constraint. It can generate turbulence-aware topology candidates, compare them with laminar surrogate designs, test objective and penalization choices, and expose whether a candidate remains robust when the turbulence closure is changed. The contribution of this thesis is therefore a reproducible bridge between turbulence-model knowledge, open-source finite-element optimization, and engineering judgment about when a turbulent topology-optimized design is credible.

\section{Limitations}
The limitations below define how far the conclusions can be taken. They are separated from the future-work actions so that the same point is not repeated twice.
\begin{enumerate}
	\item \textbf{Simplified turbulence physics and two-dimensional scope:} The framework uses a wall-resolved one-equation SA closure in 2D planar domains. This suppresses Dean vortices, secondary motion, Reynolds-stress anisotropy, and three-dimensional separation structures that can be important in physical internal-flow devices.
	\item \textbf{Frozen adjoint and sensitivities:} The finite-difference and Taylor checks verify the derivative used by MMA under the same frozen-turbulence assumption. They do not verify the derivative of a fully coupled turbulence, wall-distance, and geometry-extraction pipeline. The optimizer is therefore internally consistent with its chosen approximation, but the approximation remains a modeling decision.
	\item \textbf{Diffuse-to-sharp geometry transfer:} The optimized design is a continuous Brinkman-penalized density field, while validation is performed on an extracted sharp-wall geometry. The chosen threshold and smoothing procedure affect wall placement, local curvature, and separation behavior, so FEniCS-Ansys discrepancies cannot always be attributed to solver physics alone.
	\item \textbf{Validation level:} Independent Ansys Fluent validation with SA and SST $k-\omega$ is a necessary engineering step, but it is still numerical validation. Experimental pressure-drop and flow-field measurements would be required before using these layouts as physically validated device designs.
	\item \textbf{Optimization robustness:} The reported turbulent pipe-bend and U-bend topologies should be interpreted as optimization endpoints rather than converged final results. Although pressure drop and dissipation are physically related in steady internal flows, domain-integrated dissipation can provide a less localized optimization signal than an average inlet-pressure objective.
\end{enumerate}

\section{Future Work}
The verified turbulent fluid mechanics framework developed in this thesis establishes a foundation for several targeted extensions:

\subsection{Benchmark Completion and Optimization Robustness}
The short $70$-iteration, non-continuation pipe-bend run and the $275$-iteration turbulent U-bend endpoint reported in Chapter~\ref{ch:benchmark_results} are useful as provisional validation states, but they should be replaced by converged or deliberately polished runs before making final quantitative claims about the FEniCS optima. Future runs should use a longer MMA budget together with continuation of the RAMP, projection, move-limit, and topology-wall penalty parameters. In particular, the turbulent-design U-bend should be repeated as a more robust optimization campaign tuned for the forced $180^\circ$ return-flow geometry. The same future campaign should rerun the frozen-sensitivity diagnostic on the exact continuation settings used for the final design.

\subsection{Conjugate Heat Transfer (CHT) and Thermal Management}
The ultimate engineering application for high-Reynolds topology optimization is not merely fluid routing, but active thermal management. Historically, the immense computational cost of resolving Conjugate Heat Transfer (CHT) has forced optimization frameworks to rely on laminar approximations (e.g., natural convection as demonstrated by Alexandersen et al. \cite{Alexandersen2014}) or simplified flow proxies such as Darcy models \cite{Zhao2018}. Even when high-resolution, density-based solvers are successfully deployed for microelectronics cooling, they typically circumvent the complexities of turbulent momentum transport \cite{Barroca2023}. 

However, high-performance heat sinks still require models that capture turbulent convection. Recent work has introduced turbulence into thermal topology optimization using extended finite element (XFEM) level-set approaches with the SA model \cite{Maute2023}, computationally reduced multi-layer thermofluid approximations \cite{Zhao2021}, and standard density-based frameworks coupled to two-equation $k-\epsilon$ closures \cite{Sun2023}.

With the turbulent momentum and SA transport equations implemented and verified within the current FEniCS architecture, the framework provides a practical starting point for turbulent CHT optimization. Extending the current methodology to CHT will require implementing a turbulent Prandtl number ($Pr_t$) approach to model the turbulent heat flux vector, alongside deriving the corresponding thermal adjoint sensitivities to minimize peak temperatures or maximize overall heat transfer coefficients. Furthermore, the capacity for optimizing two-fluid heat exchangers, currently dominated by laminar assumptions \cite{Andreasen2020}, could be improved by adapting this turbulent solver.

\subsection{Three-Dimensional Topologies}
As noted in the limitations, turbulent secondary flows are inherently 3D phenomena. To fully evaluate the geometric limits of the turbulence models and design physically realizable internal-flow components, the current 2D FEniCS implementation must be extended to three dimensions. The mathematical structure of the governing equations and adjoint derivations remains similar, but the computational cost changes by orders of magnitude because wall-resolved turbulent optimization requires fine near-wall resolution over a three-dimensional design domain. Large-scale topology optimization studies have demonstrated that parallel finite-element optimization at extreme design-variable counts is possible when the formulation is paired with efficient domain decomposition and scalable solvers \cite{Aage2017}. A 3D turbulent extension of the present work should therefore prioritize parallel mesh generation, robust algebraic multigrid preconditioning, and automated sharp-geometry extraction.

\subsection{Adjoint Model Hierarchies}
This thesis deliberately uses the frozen-turbulence adjoint as the robust baseline. Future frameworks should investigate stable ways to add selected SA transport sensitivities, especially in separated flows where the constant-eddy-viscosity assumption is least accurate. Existing turbulent topology optimization studies already span several levels of adjoint complexity, from frozen or constant-eddy-viscosity approximations \cite{Dilgen2018,Zymaris2009} to more complete RANS sensitivity formulations and implementation studies \cite{Margetis2024,Wu2023}. A practical next step is therefore not necessarily to jump directly to a fully coupled turbulence adjoint, but to build an adjoint hierarchy: frozen SA sensitivities as the baseline, selected SA transport sensitivities as an intermediate model, and two-equation closures such as SST $k-\omega$ for cases where curvature, separation, and anisotropic shear dominate the design response.
